\section{Experimental Method and Evidence Handling}
\label{sec:methods}
\subsection{Saved Scenario Matrix}
The evaluation is a read-only analysis of existing experiment artifacts,
frozen for manuscript preparation on 5 September 2026. No new PSCAD
simulations were launched while producing the paper. There is one saved
deterministic result per matrix condition, so no confidence interval,
stochastic reliability estimate, or repeated-run performance distribution is
reported.

The six conditions are: S1, no-fault baseline; S2, N03 isolation at 5 s without
restoration; S3, the same isolation followed by tie closure at 6 s; S4, attempted
closure of a tie that would create a loop in the initial topology; S5, isolation
at 5 s and deliberately delayed tie closure at 45 s; and S6, a native fault
at 4.8--5.2 s with isolation at 5 s and restoration at 6 s. S2 and S3 retain
``fault'' terminology in some originating filenames or scenario descriptions,
but contain only logical isolation commands, not an applied physical fault.
Only S6 energizes the native fault component.

The input reason string in S6 refers to protection detecting a fault. This is
scenario annotation, not evidence of a simulated detector. Both fault timing
and controller timing are provided in advance in the same experiment input.
The 0.2-s interval between fault inception and isolation is imposed; it does
not measure relay decision time or communication delay.

\begin{table}[tb]
\centering\footnotesize
\caption{Saved experiment matrix. Accepted commands are shown as accepted/issued.
Events count individual breaker transitions. The N04 duration is time below
200 V; $\geq$ marks an interval still ongoing when the run ended.}
\label{tab:scenarios}
\begin{tabularx}{\textwidth}{@{}lXrrrrr@{}}
\toprule
ID & Condition & Run (s) & Cmd. & Events & Samples/ch. & N04 (s) \\
\midrule
S1 & Baseline & 1 & 0/0 & 0 & 20,001 & 0.0000 \\
S2 & Isolation only & 6 & 1/1 & 3 & 120,001 & $\geq 0.9835$ \\
S3 & Tie restoration & 7 & 2/2 & 5 & 140,001 & 0.9984 \\
S4 & Loop rejection & 6 & 0/1 & 0 & 120,001 & 0.0000 \\
S5 & Delayed restoration & 46 & 2/2 & 5 & 920,001 & 39.9984 \\
S6 & Physical fault & 7 & 2/2 & 5 & 140,001 & 1.1948
\\
\bottomrule
\end{tabularx}
\end{table}

\subsection{Measured Channels and Interruption Metric}
Each run provides six bus RMS voltage channels in kV; seven branch RMS-current
channels in kA; seven branch active-power channels in MW; seven breaker-control
state channels; three instantaneous fault-current channels in kA; and one
fault-state channel. These sum to 31 research channels. The binary
\code{State_*} channels record scheduled control signals, not independent
mechanical-contact feedback. The fault-state signal uses 1 for active,
unlike the breaker convention in Equation~\ref{eq:polarity}.

The raw output also contains animation-related traces. The existing reader
distinguishes the research recorders from sparsely sampled display channels
instead of treating every exported trace as a full-rate measurement. The saved
integration evidence reports 146 raw traces and 31 selected research traces.

For recorded bus RMS voltage $V_i[n]$, define a threshold indicator
\begin{equation}
 z_i[n]=\begin{cases}
 1,& t_n\geq 0.5\unit{s}\ \text{and}\ V_i[n]<0.2\unit{kV},\\
 0,&\text{otherwise}.
 \end{cases}
\end{equation}
An interval begins at the first below-threshold sample and ends at the first
sample no longer below threshold. The implemented analysis merges consecutive
intervals separated by at most 0.02 s, then discards intervals shorter than
0.01 s. If an interval continues to the last sample, the saved duration is
right-censored at the run end. For each retained interval $j$,
\begin{equation}
 D_{i,j}=t_{i,j}^{\mathrm{end}}-t_{i,j}^{\mathrm{start}}.
\end{equation}
Because of the merge rule, $D_{i,j}$ can include very short recovered gaps;
it is a sustained-interruption envelope, not necessarily the exact integral of
$z_i$. The threshold is 0.5 per unit of the nominal 400-V voltage, chosen for
this experiment. It is not a claim of compliance with a voltage-quality
standard, nor a customer reliability index such as SAIDI.

The peak absolute current for fault phase $\phi$ is
\begin{equation}
 I_{\phi}^{\mathrm{pk}}=\max_n |i_{\phi}[n]|.
\end{equation}
The pipeline computes peaks and threshold intervals from the full recorded
series before reducing the displayed traces. Sample alignment is
$50\unit{\mu s}$; this spacing is not a calibrated accuracy bound on RMS
filter dynamics or switching-device behavior.

\subsection{Provenance and Visualization}
The source result file stores final values, extrema, event-aligned samples,
interruption metrics, and a 501-point evenly indexed preview per channel.
The manuscript script verifies channel counts, sample counts, preview order,
finite values, final-value consistency, and interval arithmetic against this
snapshot. It does not reconstruct full-rate waveforms from those previews.
The numerical tables use stored full-rate summaries; the plotted voltage and
RMS-current trajectories use explicitly labeled decimated previews.

In particular, evenly subsampled instantaneous fault current can alias the
50-Hz waveform and miss a peak. The paper therefore presents fault-current
peak bars from the saved full-rate metrics rather than inventing a detailed
waveform between preview points. Publication-grade transient waveform figures
will require separately archived raw outputs.

Provenance includes the inspected repository revision, SHA-256 hashes of the
read input artifacts and relevant source modules, the electrical build report,
and the scenario inputs. A recorded graphical-validation artifact supplies
additional evidence for baseline selection, physical-fault selection, and a
command edit through the dashboard. This evidence is reported as an earlier
integration check, not a new GUI experiment conducted during manuscript
preparation.
